% Generated draft: requires mathematical and source review.
\documentclass[11pt]{article}
\usepackage[margin=25mm]{geometry}
\usepackage{fontspec}
\setmainfont{Latin Modern Roman}
\usepackage{amsmath,amssymb,mathtools,unicode-math}
\setmathfont{Latin Modern Math}
\usepackage{enumitem,needspace,xurl,hyperref}
\setlength{\parindent}{0pt}
\setlength{\parskip}{5pt}
\setlength{\emergencystretch}{3em}
\newcommand{\N}{\mathbb{N}}
\newcommand{\Z}{\mathbb{Z}}
\newcommand{\Q}{\mathbb{Q}}
\newcommand{\R}{\mathbb{R}}
\newcommand{\C}{\mathbb{C}}
\newcommand{\F}{\mathbb{F}}
\newcommand{\A}{\mathbb{A}}
\newcommand{\PP}{\mathbb{P}}
\newcommand{\E}{\mathbb{E}}
\newcommand{\eps}{\varepsilon}
\newcommand{\dd}{\,\mathrm{d}}
\begin{document}
\section*{UnsolvedMath ID: 11000219}
\textbf{Source problem code:} catalogue entry
\subsection{Definitions and mathematical statement}
The complete mathematical text supplied for Problem 12 is only the heading ``Analog of Ratner theorem.'' It gives no acting group, space, measure class or conclusion, so no unique theorem can be recovered from that heading alone.

The research summary identifies the intended related setting as \(SL(2,\R)\) acting on moduli of translation surfaces and attributes a Ratner-type result to Eskin--Mirzakhani--Mohammadi. In that customary framework, a translation surface is a pair \((X,\omega)\), with \(X\) a compact Riemann surface and \(\omega\) a nonzero holomorphic one-form. A stratum fixes the orders of its zeros, and one often fixes area \(\frac{i}{2}\int_X\omega\wedge\overline\omega=1\). The real-linear action on the real and imaginary parts of period coordinates defines the \(SL(2,\R)\)-action.

A precise Ratner-type question in this framework would identify the structure of orbit closures and/or ergodic invariant probability measures, for example whether they are affine in period coordinates in the appropriate immersed-orbifold sense. Affine here refers to real-linear local period equations on the corresponding non-area-normalized cone, with the associated locally Lebesgue measure; it does not mean that every such locus is a homogeneous space in the original Lie-group sense. This is a conditional interpretation suggested by the summary, not a theorem statement present in the heading.

The supplied summary reports an orbit-closure/measure-classification analogue as proved and largely answering the problem. It supplies no exact hypotheses or bibliography, and does not establish that every possible unipotent, horocycle, quantitative or other ``Ratner analogue'' is intended or settled. The missing surrounding text is needed to match the reported result to this particular source item.
\subsection{Short English statement}
The source preserves only a Ratner-analogue heading. Its summary points to the \(SL(2,\R)\) translation-surface orbit-closure and invariant-measure results of Eskin--Mirzakhani--Mohammadi, but the exact requested theorem is missing.
\subsection{Sources}
\begin{itemize}
\item Local UnsolvedMath registry, permanent ID \(11000219\), exact source code \texttt{catalogue entry}.
\item Farb, \emph{Problems on Mapping Class Groups and Related Topics}, Problem 12, ``Analog of Ratner theorem,'' PDF page 260: \url{https://www.math.uchicago.edu/~farb/papers/mcgbook.pdf}. Registry access 2026-07-29; pdf-font-structure extraction.
\item Eskin--Mirzakhani--Mohammadi are credited in the research summary with the relevant \(SL(2,\R)\) analogue. No precise paper title, theorem number or identifier is supplied.
\end{itemize}
\end{document}
